% SGA 3, Expose VII_B, English reconstruction -- Loop 1.
% Scope: Sections 5.3 and 5.4 in full on authority local pages 95--97;
% stops immediately before Section 5.5.
% Combined-reader pages: 608--610. Re-edition printed pages: 597--599.
% Original-edition marginal locator in scope: 563.
% Editorial notes in scope: 165--168.
% Authority: Polo--Gille Exp7B-13oct24.pdf,
% SHA-256 D303690E0195654BBE7C3DAA5D89EA23E24CD709A3575217234AC19AE8836694.
% Jacob Reinhold's CC BY 4.0 Markdown was consulted as comparison lineage.
% No genuine diagram occurs in this component; all mathematics is native LaTeX.

\subsection*{5.3. 注}
\phantomsection
\label{sga3:VIIB:5.3}
\addcontentsline{toc}{subsection}{5.3. 注}

\noindent\SGAInlineRefTarget{sga3:VIIB:5.3:item:a}\textup{(a)}
若一个预有限 \(k\)-代数为一族截断形式幂级数代数的完备化张量积，就称它为\emph{定常的}。\footnote{\label{sga3:VIIB:nde:165}作者无疑是指形如
\[
  \mathcal A
  =
  \mathop{\widehat{\bigotimes}}\nolimits^{c}_{n\in\mathbb N^*}
  \frac{k[[\omega'_n]]}
       {\bigl((x^{p^n})\bigr)_{x\in\omega'_n}},
\]
的完备化张量积，其中各 \(\omega'_n\) 为任意伪紧向量空间。此时 \(\omega=\omega_{\mathcal A/k}\) 与 \(\prod_n\omega'_n\) 等同，而滤过为 \(\omega_n=\prod_{i=1}^{n}\omega'_i\)；这是 \SGARefLink{sga3:VIIB:5.2.1.B}{5.2.1.B} 的情形。因此，与其称这些代数为\emph{定常的}，不如称为\emph{稳定的}；参见 \textup{[Di73]}，II，\(\S\)\SGARefLink{sga3:VIIB:2.9}{2.9}，第~75~页。例如，若对每个 \(n\in\mathbb N^*\) 都有 \(\dim\omega'_n=1\)，则 \(\mathcal A\) 表示如下函子：对 \(C\in\operatorname{Alf}/k\)，联系由 \(C\) 的元素组成的序列 \((x_n)_{n\in\mathbb N^*}\)，使得 \(x_n^{p^n}=0\)，且对充分大的 \(n\) 有 \(x_n=0\)。这一 \(\omega=\prod_{n\in\mathbb N^*}\omega'_n\) 的情形，在连通余交换 Hopf 代数的对偶背景中，对应于 M.~E.~Sweedler 研究的情形；参见 \textup{[Sw67]}，定理~3。}

设 \(G\) 为无穷小形式 \(k\)-群，\(B\) 为 \(G\) 的某个齐性空间的仿射代数。由 \SGARefLink{sga3:VIIB:5.2:theorem}{定理~5.2}，对每个 \(r\in\mathbb N\)，代数
\[
  B
  \big/
  \bigl((x^{p^r})\bigr)_{x\in\mathfrak r(B)}
\]
都是定常的。特别地，\(B\) 为定常代数的逆极限。\footnote{\label{sga3:VIIB:nde:166}这样的逆极限本身未必是上一条编者注意义下的稳定预有限 \(k\)-代数。例如，令 \(S\) 为由 \(k\) 的元素组成的序列 \((u_1,u_2,\ldots)\) 的通常 \(k\)-向量空间，并置 \(\omega=S^*\)。则 \(\omega\) 为 \(\operatorname{PC}(k)\) 中各个 \(k\) 的副本 \(kx_n\) 在 \(n\in\mathbb N^*\) 上的直和，其中 \(x_n(u)=u_n\)。因此这是 \SGARefLink{sga3:VIIB:5.2.1.A}{5.2.1.A} 的情形。代数
\[
  \mathcal A
  =
  k[[\omega]]
  \big/
  \bigl((x_n^{p^n})\bigr)_{n\in\mathbb N^*}
\]
满足 \(\omega_{\mathcal A/k}=\omega\) 和 \(\omega_n=\prod_{i=1}^{n}kx_i\)，但它不是稳定的：\(\operatorname{Spf}(\mathcal A)\) 表示如下函子，对每个 \(C\in\operatorname{Alf}/k\) 联系由 \(C\) 的元素组成的无限序列 \((x_n)_{n\in\mathbb N^*}\)，使得对每个 \(n\) 有 \(x_n^{p^n}=0\)。}

\medskip
\noindent\textup{(b)}
用 \SGARefLink{sga3:VIIB:5.2.1:section-proposition}{5.2.1} 的记号，我不知道能否选取 \(A\) 和 \(B\)，使得不存在与滤过相容的截面
\[
  \sigma:\omega\longrightarrow I
\]
\footnote{%
\label{sga3:VIIB:nde:167}除 \SGARefLink{sga3:VIIB:5.2.1.A}{5.2.1.A} 和 \SGARefLink{sga3:VIIB:5.2.1.B}{5.2.1.B} 所考虑的情形外，编者也不知道。}
不过，\(\omega\) 可以是任意赋以闭子空间递增序列所给滤过的伪紧向量空间。事实上，给定
\[
  \omega_1\subset\omega_2\subset\cdots\subset\omega,
\]
便可在
\[
  B=A
  =
  k[[\omega]]
  \big/
  \bigl((x^{p^r})\bigr)_{x\in\omega_r}
\]
上定义满足 \SGARefLink{sga3:VIIB:2.1}{2.1} 的条件 \textup{\SGARefLink{sga3:vii:paragraph:5-1:equation:i}{(i)}}、\textup{\SGARefLink{sga3:vii:paragraph:5-1:equation:ii}{(ii)}} 和 \textup{(iii)} 的对角态射
\[
  \Delta_A:A\longrightarrow A\widehat\otimes_kA
\]
只需当 \(y\) 为 \(\omega\) 中元素在 \(A\) 中的像时，置
\[
  \Delta_A(y)
  =
  y\widehat\otimes1+1\widehat\otimes y
\]

\subsection*{5.4. 推论}
\phantomsection
\label{sga3:VIIB:5.4}
\addcontentsline{toc}{subsection}{5.4. 推论}

\noindent\SGAInlineRefTarget{sga3:vii:corollary:5-4}\textbf{推论 5.4.}
{\itshape 设 \(G\) 为特征 \(p>0\) 的完美域 \(k\) 上的代数群，\(H\) 为 \(G\) 的代数子群，\(e\) 为 \(G\) 的单位元在 \(G/H\) 中的像，\(A\) 为 \(G/H\) 在 \(e\) 处的局部环。则 \(\widehat A\) 同构于形如
\[
  k[[X_1,\ldots,X_r,\ldots,X_s]]
  \big/
  (X_1^{p^{n_1}},\ldots,X_r^{p^{n_r}}).
\]
的代数。}

\medskip
\noindent\textit{证明。}
考虑无穷小形式群
\[
  \widehat G
  =
  \operatorname{Spf}(\widehat{\mathcal O}_{G,e}),
  \qquad
  \widehat H
  =
  \operatorname{Spf}(\widehat{\mathcal O}_{H,e}).
\]
由 \SGARefLink{sga3:VIIB:1.3.4}{1.3.4}，
\[
  A=\mathcal O_{G/H,e}
\]
的完备化 \(\widehat A\) 同构于
\[
  \mathscr A(\widehat G/\widehat H),
\]
推论遂由 \SGARefLink{sga3:VIIB:5.2:theorem}{定理~5.2}\textup{\SGARefLink{sga3:VIIB:5.2:theorem:item:ii}{(ii)}} 得出。\footnote{%
\label{sga3:VIIB:nde:168}另见 \textup{[DG70]}，\(\S\)III.3，定理~6.1。}
